\chapter{DF-066 Audit data visualization family across real surfaces}\label{record-107}

\noindent{\footnotesize\color{BookMuted}\path{docs/design/DF-066-audit-data-visualization-family-across-real-surfaces.md}\par}\smallskip

\begin{RecordMeta}
title: DF-066 Audit data visualization family across real surfaces
legend: DF
lane: design
priority: medium
keywords:
  - component-audit
  - dogfood
  - lowerings
  - visualization

\end{RecordMeta}

\section{Framing}\label{record-107-framing}

The data visualization family proves that Bijou can render dense numeric signals without making glyph art the only source of truth. \texttt{s\allowbreak{}p\allowbreak{}a\allowbreak{}r\allowbreak{}k\allowbreak{}l\allowbreak{}i\allowbreak{}n\allowbreak{}e\allowbreak{}(\allowbreak{})}, \texttt{b\allowbreak{}r\allowbreak{}a\allowbreak{}i\allowbreak{}l\allowbreak{}l\allowbreak{}e\allowbreak{}C\allowbreak{}h\allowbreak{}a\allowbreak{}r\allowbreak{}t\allowbreak{}S\allowbreak{}u\allowbreak{}r\allowbreak{}f\allowbreak{}a\allowbreak{}c\allowbreak{}e\allowbreak{}(\allowbreak{})}, \texttt{s\allowbreak{}t\allowbreak{}a\allowbreak{}t\allowbreak{}s\allowbreak{}P\allowbreak{}a\allowbreak{}n\allowbreak{}e\allowbreak{}l\allowbreak{}S\allowbreak{}u\allowbreak{}r\allowbreak{}f\allowbreak{}a\allowbreak{}c\allowbreak{}e\allowbreak{}(\allowbreak{})}, and \texttt{p\allowbreak{}e\allowbreak{}r\allowbreak{}f\allowbreak{}O\allowbreak{}v\allowbreak{}e\allowbreak{}r\allowbreak{}l\allowbreak{}a\allowbreak{}y\allowbreak{}S\allowbreak{}u\allowbreak{}r\allowbreak{}f\allowbreak{}a\allowbreak{}c\allowbreak{}e\allowbreak{}(\allowbreak{})} should use rich terminal geometry when available, then lower to explicit numeric summaries when visual density is unavailable.

This audit verifies the existing DOGFOOD data visualization stories against real story rendering, component-family doctrine, and mode-aware lowerings. It does not redesign charting; it turns current data-viz runtime truth into executable release proof.

\section{Sponsored Users}\label{record-107-sponsored-users}

\begin{itemize}
\tightlist
\item
  TUI app authors building dashboards, performance overlays, and status summaries from first-party visualization blocks.
\item
  Docs readers checking what happens when Unicode blocks, Braille cells, borders, and color are not meaningful output channels.
\item
  Maintainers who need one focused regression tying data-viz story metadata, variants, and component-family guidance to real output.
\end{itemize}

\section{Hills}\label{record-107-hills}

\begin{enumerate}
\def\labelenumi{\arabic{enumi}.}
\tightlist
\item
  A builder can inspect DOGFOOD and see when to use inline trends, dense area charts, metric panels, or the prebuilt performance overlay.
\item
  A reader can switch every data-viz story through interactive, static, pipe, and accessible profiles while preserving the numeric facts behind the visual shape.
\item
  A maintainer can run one cycle test that proves data-viz story metadata and design-system guidance have not drifted away from runtime behavior.
\end{enumerate}

\section{Playback Questions}\label{record-107-playback-questions}

\begin{itemize}
\tightlist
\item
  Does DOGFOOD expose the data visualization family through all four shipped stories?
\item
  Do variants cover inline sparklines, fixed ranges, Braille charts, metric panels, and performance overlays?
\item
  Do rich/static previews preserve terminal-native visual density?
\item
  Do pipe lowerings replace visual density with explicit numeric summaries?
\item
  Do accessible lowerings describe the same trend and metric facts in reading order?
\item
  Do component-family docs still position the family as informative numeric visualization rather than decoration?
\end{itemize}

\section{Requirements}\label{record-107-requirements}

\begin{itemize}
\tightlist
\item
  Keep the existing story identities: \texttt{s\allowbreak{}p\allowbreak{}a\allowbreak{}r\allowbreak{}k\allowbreak{}l\allowbreak{}i\allowbreak{}n\allowbreak{}e}, \texttt{b\allowbreak{}r\allowbreak{}a\allowbreak{}i\allowbreak{}l\allowbreak{}l\allowbreak{}e\allowbreak{}-\allowbreak{}c\allowbreak{}h\allowbreak{}a\allowbreak{}r\allowbreak{}t}, \texttt{s\allowbreak{}t\allowbreak{}a\allowbreak{}t\allowbreak{}s\allowbreak{}-\allowbreak{}p\allowbreak{}a\allowbreak{}n\allowbreak{}e\allowbreak{}l}, and \texttt{p\allowbreak{}e\allowbreak{}r\allowbreak{}f\allowbreak{}-\allowbreak{}o\allowbreak{}v\allowbreak{}e\allowbreak{}r\allowbreak{}l\allowbreak{}a\allowbreak{}y}.
\item
  Keep the shared \texttt{d\allowbreak{}a\allowbreak{}t\allowbreak{}a\allowbreak{}-\allowbreak{}v\allowbreak{}i\allowbreak{}s\allowbreak{}u\allowbreak{}a\allowbreak{}l\allowbreak{}i\allowbreak{}z\allowbreak{}a\allowbreak{}t\allowbreak{}i\allowbreak{}o\allowbreak{}n} coverage family id.
\item
  Render every data-viz story variant in every canonical story profile.
\item
  Treat block glyphs, Braille glyphs, and borders as visual-profile affordances only.
\item
  Keep pipe and accessible lowerings numeric-fact-first and semantic-preserving.
\item
  Confirm \texttt{d\allowbreak{}o\allowbreak{}c\allowbreak{}s\allowbreak{}/\allowbreak{}d\allowbreak{}e\allowbreak{}s\allowbreak{}i\allowbreak{}g\allowbreak{}n\allowbreak{}-\allowbreak{}s\allowbreak{}y\allowbreak{}s\allowbreak{}t\allowbreak{}e\allowbreak{}m\allowbreak{}/\allowbreak{}c\allowbreak{}o\allowbreak{}m\allowbreak{}p\allowbreak{}o\allowbreak{}n\allowbreak{}e\allowbreak{}n\allowbreak{}t\allowbreak{}-\allowbreak{}f\allowbreak{}a\allowbreak{}m\allowbreak{}i\allowbreak{}l\allowbreak{}i\allowbreak{}e\allowbreak{}s\allowbreak{}.\allowbreak{}m\allowbreak{}d} still matches the runtime story posture.
\end{itemize}

\section{Acceptance Criteria}\label{record-107-acceptance-criteria}

\begin{itemize}
\tightlist
\item
  \texttt{t\allowbreak{}e\allowbreak{}s\allowbreak{}t\allowbreak{}s\allowbreak{}/\allowbreak{}c\allowbreak{}y\allowbreak{}c\allowbreak{}l\allowbreak{}e\allowbreak{}s\allowbreak{}/\allowbreak{}D\allowbreak{}F\allowbreak{}-\allowbreak{}0\allowbreak{}6\allowbreak{}6\allowbreak{}/\allowbreak{}d\allowbreak{}a\allowbreak{}t\allowbreak{}a\allowbreak{}-\allowbreak{}v\allowbreak{}i\allowbreak{}s\allowbreak{}u\allowbreak{}a\allowbreak{}l\allowbreak{}i\allowbreak{}z\allowbreak{}a\allowbreak{}t\allowbreak{}i\allowbreak{}o\allowbreak{}n\allowbreak{}-\allowbreak{}f\allowbreak{}a\allowbreak{}m\allowbreak{}i\allowbreak{}l\allowbreak{}y\allowbreak{}-\allowbreak{}a\allowbreak{}u\allowbreak{}d\allowbreak{}i\allowbreak{}t\allowbreak{}.\allowbreak{}t\allowbreak{}e\allowbreak{}s\allowbreak{}t\allowbreak{}.\allowbreak{}t\allowbreak{}s} proves the cycle doc carries the modern playback sections.
\item
  The test locates all four DOGFOOD data-viz stories and verifies expected variants.
\item
  Every variant renders non-empty output in every documented profile.
\item
  Rich/static output preserves the expected visual affordance for each story.
\item
  Pipe and accessible output do not depend on block or Braille glyphs.
\item
  Pipe and accessible output preserve sample counts, ranges, latest values, trend notes, metric labels, and performance facts.
\item
  Component-family docs retain numeric trend, chart, metric panel, performance overlay, fixed-axis, pipe, and accessible guidance.
\end{itemize}

\section{Implementation Outline}\label{record-107-implementation-outline}

\begin{enumerate}
\def\labelenumi{\arabic{enumi}.}
\tightlist
\item
  Add a DF-066 cycle test that renders all data-viz DOGFOOD stories through the shared story protocol.
\item
  Read-test the data visualization section in \texttt{d\allowbreak{}o\allowbreak{}c\allowbreak{}s\allowbreak{}/\allowbreak{}d\allowbreak{}e\allowbreak{}s\allowbreak{}i\allowbreak{}g\allowbreak{}n\allowbreak{}-\allowbreak{}s\allowbreak{}y\allowbreak{}s\allowbreak{}t\allowbreak{}e\allowbreak{}m\allowbreak{}/\allowbreak{}c\allowbreak{}o\allowbreak{}m\allowbreak{}p\allowbreak{}o\allowbreak{}n\allowbreak{}e\allowbreak{}n\allowbreak{}t\allowbreak{}-\allowbreak{}f\allowbreak{}a\allowbreak{}m\allowbreak{}i\allowbreak{}l\allowbreak{}i\allowbreak{}e\allowbreak{}s\allowbreak{}.\allowbreak{}m\allowbreak{}d}.
\item
  Move the backlog note into \texttt{d\allowbreak{}o\allowbreak{}c\allowbreak{}s\allowbreak{}/\allowbreak{}d\allowbreak{}e\allowbreak{}s\allowbreak{}i\allowbreak{}g\allowbreak{}n\allowbreak{}/} and update the v6 lane pointer.
\item
  Record any drift as follow-on backlog debt instead of widening this audit.
\end{enumerate}

\section{Drift Check}\label{record-107-drift-check}

No runtime drift was found in this slice. The existing data visualization stories already render visual glyph density in interactive/static profiles, lower sparkline and Braille previews to numeric summaries in pipe mode, and linearize metric and performance summaries in accessible mode.

No follow-on backlog item is needed for this slice.

\section{Playback}\label{record-107-playback}

\begin{itemize}
\tightlist
\item
  RED: the release lane had only a backlog note; no DF-066 playback test tied all shipped data-viz stories to every canonical profile.
\item
  GREEN: the new cycle test renders all \texttt{s\allowbreak{}p\allowbreak{}a\allowbreak{}r\allowbreak{}k\allowbreak{}l\allowbreak{}i\allowbreak{}n\allowbreak{}e}, \texttt{b\allowbreak{}r\allowbreak{}a\allowbreak{}i\allowbreak{}l\allowbreak{}l\allowbreak{}e\allowbreak{}-\allowbreak{}c\allowbreak{}h\allowbreak{}a\allowbreak{}r\allowbreak{}t}, \texttt{s\allowbreak{}t\allowbreak{}a\allowbreak{}t\allowbreak{}s\allowbreak{}-\allowbreak{}p\allowbreak{}a\allowbreak{}n\allowbreak{}e\allowbreak{}l}, and \texttt{p\allowbreak{}e\allowbreak{}r\allowbreak{}f\allowbreak{}-\allowbreak{}o\allowbreak{}v\allowbreak{}e\allowbreak{}r\allowbreak{}l\allowbreak{}a\allowbreak{}y} variants across interactive, static, pipe, and accessible profiles.
\item
  Rich/static previews preserve block, Braille, or boxed metric affordances as appropriate.
\item
  Pipe previews preserve numeric summaries without block or Braille glyphs.
\item
  Accessible previews linearize the same metric and trend facts in reading order.
\end{itemize}

\section{Retrospective}\label{record-107-retrospective}

The visualization components did not need runtime changes. The useful release work was proving that visual density is backed by data facts, so future block and layout work can reuse these components without treating charts as opaque terminal art.

